% Irodov Part V §5.5, original problems 216--223.
% Supplied Arihant PDF pp.231--232 (printed pp.222--223).
% The source uses chi for absorption and mu for narrow-beam total attenuation.
\dispersionproblem{216}{222}
单色平面光束垂直照射厚度为 $l$ 的平行平板，
材料的线吸收系数沿法线方向由入射面的 $\chi_1$
线性变到出射面的 $\chi_2$。
两面的强度反射率均为 $\rho$，忽略二次反射。
求整块平板的能流透射系数。
\exsolution
设距离入射面深度为 $z$，吸收系数为
\[
 \chi(z)=\chi_1+\frac{\chi_2-\chi_1}{l}z.
\]
比尔--朗伯定律对非均匀材料取微分形式
$dI/dz=-\chi(z)I$，积分得
\[
 I(l)=I(0)\exp\left[-\int_0^l\chi(z)\,dz\right].
\]
由于线性函数在区间上的平均值恰为
$(\chi_1+\chi_2)/2$，
加入两个界面反射损耗，最终
\[
 \boxed{T=(1-\rho)^2
 \exp\left[-\frac{\chi_1+\chi_2}{2}l\right]}.
\]
\excheck 当 $\chi_1=\chi_2=\chi$ 时，
正确退化为 $(1-\rho)^2e^{-\chi l}$；
指数无量纲。

\dispersionproblem{217}{223}
一束光以强度 $I_0$ 正入射厚度为 $l$ 的
平行平板，入射光在区间
$\lambda_1\le\lambda\le\lambda_2$ 内具有\emph{均匀谱强度}，
区间外无光。
线吸收系数 $\chi(\lambda)$ 随波长线性变化，
在两端分别为 $\chi_1,\chi_2$。
两面反射率均为 $\rho$，忽略二次反射。
求透射光的\emph{总}强度。
\exsolution
令 $\Delta\lambda=\lambda_2-\lambda_1>0$，
均匀谱强度为 $I_0/\Delta\lambda$。
每个波长分量单程透射比例
$(1-\rho)^2e^{-\chi(\lambda)l}$，因而
\[
 I_{\rm out}=\frac{I_0(1-\rho)^2}{\Delta\lambda}
 \int_{\lambda_1}^{\lambda_2}
 e^{-\chi(\lambda)l}\,d\lambda.
\]
线性变换
$u=\chi_1+(\chi_2-\chi_1)
(\lambda-\lambda_1)/\Delta\lambda$
把积分化成指数积分，得
\[
 \boxed{I_{\rm out}
 =I_0(1-\rho)^2
 \frac{e^{-\chi_1l}-e^{-\chi_2l}}
 {l(\chi_2-\chi_1)}}.
\]
若 $\chi_2=\chi_1=\chi$，
取极限而非直接除零，
\[
 \boxed{I_{\rm out}
 =I_0(1-\rho)^2 e^{-\chi l}}.
\]
\excheck 由积分平均值定理，透射比例必在
$(1-\rho)^2e^{-\max(\chi_1,\chi_2)l}$ 与
$(1-\rho)^2e^{-\min(\chi_1,\chi_2)l}$ 之间；
分子分母变号同步，不会产生负光强。

\dispersionproblem{218}{223}
一块光学滤光片厚度为 $d$，
其线吸收系数随波长满足
\[
 \chi(\lambda)=\alpha\left(1-\frac\lambda{\lambda_0}\right)^2
 \quad(\mathrm{cm}^{-1}),
\]
其中 $\alpha$ 的数值以 $\mathrm{cm}^{-1}$ 为单位。
界面反射率与波长无关。
求以 $\lambda_0$ 为中心、边缘强度衰减为中心
$\eta$ 倍（$\eta>1$）的透射带宽 $\Delta\lambda$。
\exsolution
滤光片透射强度
$I(\lambda)=C I_{\rm in}e^{-\chi(\lambda)d}$，
其中 $C$ 为同一表面反射的波长无关因子。
中心波长处 $\chi(\lambda_0)=0$，故
$I(\lambda_0)=CI_{\rm in}$。
带宽两侧定义
$I(\lambda_\pm)=I(\lambda_0)/\eta$，所以
\[
 \alpha d
 \left(1-\frac{\lambda_\pm}{\lambda_0}\right)^2
 =\ln\eta.
\]
解得
$\lambda_\pm=\lambda_0[1\pm
\sqrt{\ln\eta/(\alpha d)}]$，
两边缘之差为
\[
 \boxed{\Delta\lambda=
 2\lambda_0\sqrt{\frac{\ln\eta}{\alpha d}}}.
\]
\excheck $\alpha d$ 与 $\ln\eta$ 无量纲；
只有 $\lambda_->0$ 时，两侧边缘均落在物理正波长范围。

\dispersionproblem{219}{223}
点状单色光源位于一球壳形吸收介质的球心，
总发光功率为 $\Phi$；
球壳内半径 $a$、外半径 $b$，其中 $b>a$。
介质线吸收系数为 $\chi$，
两边界反射率都为 $\rho$。
忽略二次反射，求穿过球壳之后距球心 $r>b$
处的光强。
\exsolution
由球对称性，源发出总功率 $\Phi$，
每单位立体角功率为 $\Phi/(4\pi)$。
穿过球壳的任意径向光线在材料内经过距离
$b-a$，遭遇两个界面透射因子 $(1-\rho)^2$，
吸收因子为 $e^{-\chi(b-a)}$。
出射总功率
$\Phi_{\rm out}=\Phi(1-\rho)^2e^{-\chi(b-a)}$，
故以光线法向的单位面积上能流定义强度，
\[
 \boxed{I(r)=\frac{\Phi}{4\pi r^2}
 (1-\rho)^2e^{-\chi(b-a)}},\qquad r>b.
\]
刚离开球壳表面取 $r=b$，即
$\boxed{I(b)=\Phi(1-\rho)^2e^{-\chi(b-a)}/(4\pi b^2)}$。
\excheck 射线束总透射功率只与壳厚 $b-a$ 有关，
球面上的单位面积光强仍按 $1/r^2$ 稀释。
本题用光强表示球面\emph{辐照度}，而非点光源的光度学发光强度。

\dispersionproblem{220}{223}
波长为 $20\,\mathrm{pm}$ 的窄 X 射线束穿过厚度
$d=1.0\,\mathrm{mm}$ 的铅板，
其质量衰减系数为
$\mu/\rho=3.6\,\mathrm{cm^2/g}$。
求光束强度降低的倍数。
\exsolution
窄束服从指数衰减
$I=I_0e^{-\mu d}$，
其中 $\mu=(\mu/\rho)\rho_{\mathrm{Pb}}$。
取铅的密度约
$\rho_{\mathrm{Pb}}=11.34\,\mathrm{g/cm^3}$，
$d=0.10\,\mathrm{cm}$，得
\[
 \mu d=
 (3.6\,\mathrm{cm^2/g})
 (11.34\,\mathrm{g/cm^3})
 (0.10\,\mathrm{cm})
 =4.0824.
\]
所以
\[
 \boxed{\frac{I_0}{I}=e^{4.0824}\simeq59.3}.
\]
\excheck 指数是无量纲，且“降低倍数”应为 $I_0/I>1$；
窄束中散射到束外的光子也属于衰减。

\dispersionproblem{221}{223}
波长为 $62\,\mathrm{pm}$ 的窄 X 射线经过
厚度 $d_{\mathrm{Al}}=2.6\,\mathrm{cm}$ 的铝屏。
求为产生相同强度衰减而需要的铅屏厚度。
两材料的质量衰减系数分别为
$(\mu/\rho)_{\mathrm{Al}}=3.48$ 和
$(\mu/\rho)_{\mathrm{Pb}}=72.0\,\mathrm{cm^2/g}$。
\exsolution
若两种屏对窄束有相同的透射比例，
由比尔定律指数必须相等：
\[
 \left(\frac{\mu}{\rho}\right)_{\rm Al}
 \rho_{\rm Al}d_{\rm Al}
 =
 \left(\frac{\mu}{\rho}\right)_{\rm Pb}
 \rho_{\rm Pb}d_{\rm Pb}.
\]
因此
\[
 \boxed{d_{\rm Pb}=d_{\rm Al}
 \frac{(\mu/\rho)_{\rm Al}\rho_{\rm Al}}
 {(\mu/\rho)_{\rm Pb}\rho_{\rm Pb}}}.
\]
取密度
$\rho_{\rm Al}=2.70\,\mathrm{g/cm^3}$、
$\rho_{\rm Pb}=11.34\,\mathrm{g/cm^3}$，
\[
 \boxed{d_{\rm Pb}\simeq
 2.6\,\mathrm{cm}\times
 \frac{3.48\times2.70}{72.0\times11.34}
 \simeq0.0299\,\mathrm{cm}=0.30\,\mathrm{mm}}.
\]
\excheck 两材料厚度不同，但质量衰减系数与密度、
厚度的乘积相同，保证透射比相同。

\dispersionproblem{222}{223}
用窄单色 X 射线照射一块铝板。
若对应质量衰减系数为
$\mu/\rho=0.32\,\mathrm{cm^2/g}$，
求使射线强度减少至原来一半的铝层厚度。
\exsolution
半值层厚度 $d_{1/2}$ 满足
$e^{-\mu d_{1/2}}=1/2$，
因而
\[
 d_{1/2}=\frac{\ln2}{\mu}
 =\frac{\ln2}{(\mu/\rho)\rho_{\rm Al}}.
\]
取铝密度
$2.70\,\mathrm{g/cm^3}$，得到
\[
 \boxed{d_{1/2}=
 \frac{0.6931}{0.32\times2.70}\,\mathrm{cm}
 \simeq0.802\,\mathrm{cm}=8.0\,\mathrm{mm}}.
\]
\excheck “减少一半”对应透射比 $1/2$，指数为 $\ln2$；
不能用 $\mu d=1$ 代替半值条件。

\dispersionproblem{223}{223}
一块材料可使窄 X 射线强度降低
$\eta=50$ 倍。
问其中相当于多少个将强度减半的“半值层”？
\exsolution
设材料厚度相当于 $N$ 个半值层，
每层的透射比例为 $1/2$，
故
$I/I_0=(1/2)^N$。
题给 $I/I_0=1/\eta$，因此
\[
 2^N=\eta,\qquad
 \boxed{N=\frac{\ln\eta}{\ln2}
 =\log_2 50\simeq5.64}.
\]
\excheck 半值层的数量是等效的连续厚度比，
不必为整数；确实有
$2^{5.64}\simeq50$。
